%% General CDCL construction, with source manifestations kept distinct.
\chapter{CDCL, First-UIP Analysis, and Watched-Literal Propagation}
\label{sec:p12-cdcl-uip-and-watched-literals}

This chapter supplies the operational learning and propagation algorithms that
the apparatus in Definition~\ref{def:p12-cdcl-apparatus} does not by itself
prove.  The older handbook manuscript, the Chaff author paper, and MiniSat's
explicitly extended version~1.2 are separately attributed
\cite{MarquesSilvaLynceMalik2008CDCLManuscript,MoskewiczEtAl2001ChaffV56,EenSorenssonMiniSatExtended12}.
The older manuscript prints 2008 and pages 127--149; it is not silently
identified with the final 2009 chapter or its much longer 2021 successor.
The finite-CNF algorithm and complete proofs below state their own operational
side conditions, rather than crediting every solver bearing these names with
the same algorithm or liveness policy.

\section{Source notation and the operational state}

\begin{hbdefinition}{def:p12-cdcl-uip-state}{Formula, assignment, trail, and reasons}
Retain the older handbook's notation: \(X\) is a finite variable set,
\(\varphi\) a CNF clause set, \(\omega\) a clause, and
\(\nu:X\to\{0,u,1\}\) a partial assignment, with \(u\) meaning
unassigned, not a third propositional truth value.  For a signed literal
\(p\), \(\neg p\) denotes sign-flipping literal complement, so the
complement of \(\neg x\) is \(x\), not a new double-negation formula.
For literal evaluation,
\((\neg x)^\nu=1-\nu(x)\) with \(1-u=u\); clauses are false when all
literals are false, true when some literal is true, and unit when exactly one
literal is unassigned and every other literal false
\cite[Sec.~4.2]{MarquesSilvaLynceMalik2008CDCLManuscript}.  An empty clause
is false; an empty database is true.  Clauses are sets, so duplicate literal
occurrences are removed.  Tautological clauses can be discarded without
changing any model; if retained they cannot force or conflict under a
consistent assignment.  The proved loop retains its input clauses; optional
tautology removal is a separate model-equivalent initial normalization.

The operational state has an immutable input \(\varphi_0\), an active
database \(\varphi\) containing the input and retained learned clauses, a
consistent ordered trail \(T\), a current decision level \(d\), levels
\(\delta(x)\), reasons \(\alpha(x)\), and a queue \(Q\) of assigned
literals whose watch events remain to be processed.  A trail entry is the
literal made true; no variable occurs twice.  Level zero contains the root
propagations.  Each positive level begins with one decision and continues
with its propagations.  Set \(\delta(x)=-1\) for an unassigned variable.
\(\alpha(x)=\mathrm{NIL}\) for a decision or an
unassigned variable.  For a propagated literal \(p\), its reason clause
contains \(p\), and all other reason literals were false before \(p\)
was appended.  Reasons are retained while their trail entries exist.

The manuscript assigns a propagated variable the maximum level of the other
variables in its reason, or zero for a unit reason.  MiniSat's extended
implementation instead records the current decision level at enqueueing
\cite[Eq.~(4.6)]{MarquesSilvaLynceMalik2008CDCLManuscript}
\cite[Fig.~9]{EenSorenssonMiniSatExtended12}.  The operational core here uses
the latter, block-ordered trail convention; it always saturates propagation
before making a new decision.  These level policies are not identified for
arbitrary out-of-order propagation or incremental constraint insertion.
\index{CDCL!operational state}\index{reason clause}\index{trail}
\end{hbdefinition}

\begin{hbmethod}{def:p12-cdcl-uip-enqueue}{Enqueue, propagate, decide, and cancel}
To enqueue \(p\) with a reason \(\omega\), do nothing if \(p\) is
already true; report a conflict if it is false; otherwise set it true, append
it to \(T\) and \(Q\), and record \(\delta\) and \(\alpha\).
A decision enqueues a literal of an unassigned variable with no reason,
after incrementing \(d\), and is permitted only when \(Q\) is empty and
the database is unit-saturated and nonconflicting.  Propagation repeatedly
uses an actual unit clause, stopping at a falsified clause or a fixed point.
The watched implementation is specified below, not assumed correct by name.

Cancel to level \(b\) by erasing every trail entry with level greater than
\(b\), resetting its assignment and reason, and setting \(d=b\).
Unchanged reasons refer only to earlier retained entries.  On a conflict all
pending forward events belong to the current level; after backjumping below
that level they can be discarded, followed by enqueueing the new assertion.
A voluntary restart is made at a propagation boundary and cancels to level
zero while retaining the chosen learned clauses.  It does not flip root facts
or treat a conflict as SAT.  These are operations on a general finite input,
not a replay of a single example.
\end{hbmethod}

\begin{hblemma}{thm:p12-cdcl-uip-local-sound}{proved=annex:p12-cdcl-uip-local-sound}{Reason and resolution invariants}
Unit propagation terminates on a fixed finite database and preserves all
models extending the preceding trail.  Every reason's nonpropagated literals
occur earlier on the trail.  If two clauses are consequences of
\(\varphi_0\), their propositional resolvent on a complementary pivot pair
is also a consequence.  Thus learning by the stated sequence of resolution
steps preserves exactly the input models.  The claim concerns entailed CDCL
clauses; it does not apply to non-entailed RAT additions.
\end{hblemma}

\section{A terminating first-UIP construction}

\begin{hbdefinition}{def:p12-cdcl-uip-graph}{Implication graph and the first UIP}
The implication graph has a node for each true trail literal and a conflict
node \(\kappa\).  A propagated literal \(p\) has incoming edges from
the true complements of the other literals in its reason.  The conflict node
has incoming edges from the true complements of its falsified clause.
Edges follow trail order, so the graph is acyclic.  A current-level node is a
\emph{unique implication point} (UIP) if every path from the current
decision node to \(\kappa\) passes through it.  The \emph{first UIP} is
the downstream-most such node: no other UIP lies strictly between it and
the conflict.  This is a domination relation, not a shortest-edge-distance
definition.  Only conflict-relevant ancestors are needed for analysis.
\index{unique implication point}\index{first UIP}\index{implication graph}
\end{hbdefinition}

\begin{hbmethod}{def:p12-cdcl-uip-analyze}{Reverse-trail first-UIP analysis}
At a conflict above level zero, initialize the temporary clause \(\omega_L\)
to the falsified conflict clause.  While it has more than one current-level
literal, choose the literal \(\neg p\) whose true trail complement \(p\)
is latest among its current-level entries.  Resolve with the reason
\(\alpha(p)\):
\[
 \omega_L\ \gets
 (\omega_L\setminus\{\neg p\})\cup(\alpha(p)\setminus\{p\}).
\]
Here \(\alpha(p)\) means the reason of the variable of \(p\).
Stop as soon as exactly one current-level literal \(a\) remains.  Set
\[
 b=\max\bigl(\{0\}\cup
       \{\delta(l):l\in\omega_L\setminus\{a\}\}\bigr).
\]
Erase levels greater than \(b\), insert \(\omega_L\), and enqueue
\(a\) with that clause as its reason.  The maximum is zero for a learned
unit.  At a root conflict return UNSAT instead of trying to select a positive
level pivot.  A literal's level is its variable's level, regardless of its
polarity.  The original source's selected-resolution and UIP-count equations
are in §4.4; the reverse-trail pivot choice and assertion placement are explicit
in MiniSat's analyze/record methods
\cite[Eqs.~(4.16)--(4.18)]{MarquesSilvaLynceMalik2008CDCLManuscript}
\cite[Figs.~10--11]{EenSorenssonMiniSatExtended12}.
\end{hbmethod}

\begin{hbtheorem}{thm:p12-cdcl-uip-analysis}{proved=annex:p12-cdcl-uip-analysis}{First-UIP analysis, assertion, and backjumping}
For a conflict reached by the stated propagate-before-decide discipline at
level \(d>0\), the construction terminates.  Each intermediate clause is
entailed by the input and falsified by the conflict trail, and is therefore
non-tautological.  Every required reason exists.  The final clause has exactly
one current-level literal; its complement is the first UIP.  The computed
\(b<d\) is the lowest retained level at which the final clause is unit.
After cancellation it asserts \(a\).  The argument covers root literals and
learned units without deleting them from the intermediate clauses.
\end{hbtheorem}

\begin{hbexample}{def:p12-cdcl-uip-source-trace}{The source conflict and its two learning boundaries}
Retain the manuscript's variables and clauses:
\[
\begin{array}{lll}
 \omega_1=\{x_1,x_{31},\neg x_2\},&
 \omega_2=\{x_1,\neg x_3\},&
 \omega_3=\{x_2,x_3,x_4\},\\
 \omega_4=\{\neg x_4,\neg x_5\},&
 \omega_5=\{x_{21},\neg x_4,\neg x_6\},&
 \omega_6=\{x_5,x_6\}.
\end{array}
\]
The decisions \(x_{21}=0@2\), \(x_{31}=0@3\), and \(x_1=0@5\)
give the ordered implications \(\neg x_2\) by \(\omega_1\),
\(\neg x_3\) by \(\omega_2\), \(x_4\) by \(\omega_3\),
\(\neg x_5\) by \(\omega_4\), and \(\neg x_6\) by
\(\omega_5\).  All five are at level 5.  Earlier levels not affecting
these clauses may contain other decisions; the displayed level numbers are
the source's, not a newly contiguous three-decision renumbering.
Clause \(\omega_6\) is false.  Reverse-trail analysis first resolves it
with \(\omega_5\), producing \(\{x_5,x_{21},\neg x_4\}\), then
with \(\omega_4\), producing \(\{x_{21},\neg x_4\}\).
Only \(\neg x_4\) is at the current level, so \(x_4\) is the first UIP.
Backjumping to level 2 retains \(\neg x_{21}\), erases level-5 assignments,
and the learned clause asserts \(\neg x_4\).

The source's other admissible initial resolution order gives
\(\{x_5,x_6\}\), \(\{\neg x_4,x_6\}\), then
\(\{\neg x_4,x_{21}\}\), the same first-UIP clause.  Continuing
past the UIP, rather than stopping there, gives
\(\{x_2,x_3,x_{21}\}\), \(\{x_1,x_{31},x_3,x_{21}\}\), and
\(\{x_1,x_{31},x_{21}\}\); its backjump level is 3.
The source's all-UIP/decision-boundary learning and the first-UIP stopping
boundary are thereby kept distinct
\cite[Examples~4.2.4, 4.4.1--4.4.3; Tables~4.1--4.2]{MarquesSilvaLynceMalik2008CDCLManuscript}.
The later backjump makes \(x_4\) false, which immediately satisfies
\(\omega_4,\omega_5\); the conflict is not left on the trail.
\end{hbexample}

\section{The full two-watched-literal update}

\begin{hbdefinition}{def:p12-cdcl-uip-watches}{Watch lists, units, and attachment}
For each nonunit active clause, select two distinct literal positions,
\(w_1(\omega),w_2(\omega)\).  A list \(W(l)\) contains every clause
currently watching the literal \(l\); making \(l\) false triggers that
list.  This is the Chaff literal-oriented convention.  MiniSat instead indexes
the lists by the newly true complementary literal; the event is the same,
not a different semantic condition
\cite[Sec.~2]{MoskewiczEtAl2001ChaffV56}
\cite[Figs.~7--9]{EenSorenssonMiniSatExtended12}.
Initialize input watches before decisions, handle an empty input clause as
an immediate conflict, and enqueue input units separately at level zero.
Do not simulate a unit by two duplicate watch positions.  For a newly learned
nonunit clause, watch its asserting literal and a nonasserting literal of
maximal level \(b\), attach both list references, and immediately enqueue
the asserting literal.  Learned units are root facts and are retained for the
plain-SAT core.  The maximal-level second watch is essential to the rollback
proof, not a heuristic choice without correctness content.
\end{hbdefinition}

\begin{hbmethod}{def:p12-cdcl-uip-watch-update}{Process a false-watch event}
When dequeuing a true literal \(p\), inspect the current list
\(W(\neg p)\).  For each live clause, designate the watched false literal
\(\neg p\) and its other watch \(q\):
\begin{enumerate}
 \item If \(q\) is true, leave the watches and clause in the current list.
 \item Otherwise search all nonwatched clause positions for a literal that
       is not false.  If one exists, move the false watch to that position,
       remove the old list reference, insert the new one, and continue.
 \item If none exists, leave the false watch registered.  If \(q\) is
       unassigned, enqueue \(q\) with this clause as reason.  If \(q\)
       is false, report this clause as a conflict.
\end{enumerate}
Thus a true replacement may be chosen, not just an unassigned one.  A binary
clause has no nonwatched position and goes directly to the third case.
Queue a newly assigned literal only once; a repeated true enqueue is a no-op
and an opposite enqueue reports conflict.  When a list is temporarily removed
for traversal, restore all unprocessed references before returning a conflict;
they must not disappear because the scan was interrupted.  Detach every watch
reference when an unlocked learned clause is deleted, or retain an explicit
inactive mark tested before access.  The mathematical procedure assumes exact
live membership, not dangling references.
\end{hbmethod}

\begin{hbtheorem}{thm:p12-cdcl-uip-watch-correct}{proved=annex:p12-cdcl-uip-watch-correct}{Watched propagation and level-prefix backtracking}
Under the initialization, attachment, queue and level-block conditions just
stated, watched propagation terminates and performs exactly sound unit
inferences.  If it finishes without conflict, no active clause is false or
unit with an unassigned literal: it has reached the Boolean-propagation fixed
point.  It detects every unit/conflict event required to reach that fixed
point, though the order and chosen reasons may differ from a full scan.
After cancellation of entire decision levels, watch positions need not be
restored or moved.  This rollback claim is not for arbitrary assignment
editing, arbitrary partial erasure within a level, or improperly attached
clauses.
\end{hbtheorem}

\begin{hbexample}{def:p12-cdcl-uip-watch-failure}{Why the second learned watch cannot be arbitrary}
In this example \(u,p,q,r,s\) are propositional variables; \(u\) is not
the earlier unassigned marker.  Consider input clauses
\(\{\neg u,\neg p,r,s\}\),
\(\{\neg u,\neg p,r,\neg s\}\), and \(\{\neg r,\neg q\}\).
Decide \(p\) at level 1, \(q\) at level 2, obtaining \(\neg r\) at
level 2, then decide \(u\) at level 3.  The first clause forces \(s\)
and the second conflicts.  Resolving them on \(s\) learns
\(\omega_L=\{\neg u,\neg p,r\}\), with first UIP \(u\) and
backjump level 2.  The assertion is \(\neg u\).

An incorrect attachment watches \(\neg u\) and the lower-level literal
\(\neg p\), instead of the required level-2 literal \(r\).
At the asserting state \(\neg u\) is true at level 2 while \(\neg p\)
is false at level 1.  To test this attachment's claimed rollback property,
cancel the level prefix to level 1, erasing \(\neg u,\neg r,q\) and
retaining \(p\).  This is a rollback boundary check, not a claim that the
three-clause search would spontaneously choose that cancellation.
Deciding \(q\) again
forces \(\neg r\), making the learned clause unit in \(\neg u\).
No watched literal has newly become false: \(r\) is unwatched and
\(\neg p\) was already false.  The defective data structure misses the
required propagation.  With the correct second watch \(r\), its new
falsification triggers the clause and enqueues \(\neg u\).  The source's
constructor requirement is therefore a real side condition
\cite[Fig.~8]{EenSorenssonMiniSatExtended12}; no claim that arbitrary
backtracking makes any pair of watches safe is made.
\end{hbexample}

\section{A complete CDCL search and its actual liveness hypotheses}

\begin{hbmethod}{def:p12-cdcl-uip-search}{General finite-CNF CDCL loop}
Initialize root units and watches; run propagation.  Repeat:
\begin{enumerate}
 \item If there is a root conflict, return UNSAT.
 \item If there is a positive-level conflict, perform the complete first-UIP
       construction, cancel to its backjump level, attach the learned clause,
       enqueue its assertion, and return to propagation.
 \item If propagation is complete and all input variables are assigned,
       return the total assignment as SAT.
 \item Otherwise choose any unassigned variable and either value, make its
       next-level decision, and return to propagation.
\end{enumerate}
Unlocked learned clauses may be forgotten and voluntary restarts may occur
only under a separately declared policy.  The input clauses are retained.
A clause currently serving as a reason is locked and may not be deleted until
that reason is no longer needed, or until a correct replacement reason is
installed.  Do not return SAT merely because a lazy status cache labels every
clause unresolved rather than conflicting.  The explicit loop always
propagates after a learned assertion, including after repeated conflicts.
\end{hbmethod}

\begin{hbtheorem}{thm:p12-cdcl-uip-search-correct}{proved=annex:p12-cdcl-uip-search-correct}{CDCL soundness and a terminating complete policy}
Every SAT return of this loop satisfies the immutable input.  Every UNSAT
return is justified by an entailed root refutation.  Learning, unlocked
learned-clause deletion and restarts preserve the input model set, subject to
the retained-reason and propagation conditions.

If forgetting and voluntary restarts are disabled, this finite-CNF loop
terminates for every choice of unassigned decision variable and polarity, and
returns SAT exactly for satisfiable inputs.  More generally the same conclusion
holds with retained learning and restarts only after completed nonroot
conflicts, with no sequence of administrative steps delaying propagation or
decisions indefinitely.  Each such conflict adds a previously absent
non-tautological clause; there are only finitely many over the fixed variables.
The claim is not termination under arbitrary restarting or forgetting.
\end{hbtheorem}

For a concrete failed liveness policy, take \(\{p,q\}\), choose \(p\),
then voluntarily restart before choosing \(q\), and repeat forever.  Every
operation is semantically sound, the input is SAT, no conflict or learning
occurs, and the solver never returns a total model.  Another policy that
forgets all learned information and repeats the same bounded search can lose
the finite-clause progress argument.  A timeout is not UNSAT.

\section{Branching, look-ahead, activity, and deletion are different mechanisms}

\begin{hbdefinition}{def:p12-cdcl-uip-activity}{The separately attributed activity updates}
Chaff's 2001 VSIDS has a counter for each literal polarity, initially zero;
adding a clause increments each of its literal counters, decisions select a
highest-counter unassigned variable/polarity, default ties are randomized, and
periodically all counters are divided by a fixed constant
\cite[Sec.~3]{MoskewiczEtAl2001ChaffV56}.  The older handbook's §4.5.3
describes the corresponding clause/literal occurrence feedback.  This is not
MiniSat's polarity-blind variable activity.

In MiniSat extended~1.2, variables appearing in a recorded conflict clause are
bumped, and a learned clause used in conflict analysis has its clause activity
bumped.  Variable choice orders the unassigned variables by activity.  The
implementation stores a bump increment \(\Delta\), increases it by the
inverse decay factor after conflicts, and occasionally rescales both the
scores and \(\Delta\) together; Fig.~14 uses threshold \(10^{100}\) and
rescale factor \(10^{-100}\)
\cite[Secs.~3, 4.6; Figs.~8, 14]{EenSorenssonMiniSatExtended12}.
All these scores stand as search priorities, not truth values, clause
entailments, or certificates.  Their numerical implementations require their
own overflow/rounding and tie policies.
\end{hbdefinition}

\begin{hblemma}{thm:p12-cdcl-uip-activity-scale}{proved=annex:p12-cdcl-uip-activity-scale}{Ideal-real lazy decay preserves priority order}
Fix \(0<\rho<1\) and the same sequence of bumps.  In exact real arithmetic,
eagerly multiplying all scores by \(\rho\) after a conflict is equivalent
for score ordering to leaving scores unchanged and dividing the next bump
increment by \(\rho\), with simultaneous positive rescaling of every score
and the increment.  This preserves rankings and exact ties, not a theorem
about floating-point executions or solver performance.
\end{hblemma}

\begin{hbmethod}{def:p12-cdcl-uip-lookahead}{Look-ahead without confusing scores with inferences}
For an unassigned variable \(x\), temporarily propagate each of \(x\)
and \(\neg x\) from copies of the current assignment and database.
Discard each temporary trail afterward; do not leak its decisions into the
permanent state.  A branch conflict licenses the opposite assignment under
the current retained prefix.  If both branches imply the same literal, that
literal is implied under the prefix.  Counts of reductions or propagated
literals can also be used to score a branch, but such a score licenses no
logical assignment.  The method is the Boolean failed-literal/common-
consequence core, not a reconstruction of the entire separately allocated
2021 look-ahead or branching-heuristics chapters.
\end{hbmethod}

\begin{hblemma}{thm:p12-cdcl-uip-lookahead-sound}{proved=annex:p12-cdcl-uip-lookahead-sound}{Failed-literal and two-branch consequence soundness}
For any consistent retained prefix, a unit conflict under temporary \(x\)
shows that every extending model makes \(\neg x\) true.  If conflicts
occur under both polarities, no model extends the prefix.  If both branches
force the same literal \(l\), every extending model makes \(l\) true.
These conclusions are relative to the prefix; a globally learned clause must
include the negations of the prefix assumptions, unless a separate derivation
discharges them.
\end{hblemma}

Restart schedules decide when to abandon a trail; activity decides which
unassigned variable to choose; phase saving chooses a value; deletion decides
which entailed clauses to retain.  They are not interchangeable.  Chaff's
paper describes scheduled relevance-based lazy deletion and retained-clause
restarts; MiniSat extended~1.2 uses clause activity and excludes locked reasons
from removal.  The older handbook distinguishes bounded-size, relevance and
activity/age policies \cite[Secs.~4.1--4.2]{MoskewiczEtAl2001ChaffV56}
\cite[Secs.~4.5.2--4.5.4]{MarquesSilvaLynceMalik2008CDCLManuscript}
\cite[Secs.~4.5, 4.7--4.8]{EenSorenssonMiniSatExtended12}.
Those source-specific policies are not the no-forgetting theorem's hypotheses.
Likewise the extended MiniSat constraint and incremental-assumption interfaces
require reason-soundness and assumption-relative guarantees beyond the plain
CNF loop proved here; they are not silently proved by a shared solver name.
